% 相同边缘Bernoulli(1/2)的三种联合分布（教学数据，可在Omega=[0,1)上同时实现）：
% (a) 同向 Y=X：p(0,0)=p(1,1)=1/2；P(X=Y)=1，Cov=1/4。
% (b) 反向 Y=1-X：p(0,1)=p(1,0)=1/2；P(X=Y)=0，Cov=-1/4。
% (c) 独立：四格各1/4；P(X=Y)=1/2，Cov=0。三表的行和、列和都是(1/2,1/2)。
\begin{tikzpicture}[line width=.6pt,font=\fontsize{8}{10}\selectfont,>=Stealth,
  cell/.style={draw=FigureStroke,minimum width=9mm,minimum height=6.5mm,inner sep=0pt},
  hot/.style={cell,fill=FigureBlueFill},
  marg/.style={cell,fill=FigureGreenFill,minimum width=8mm}]
\foreach \dx/\a/\b/\c/\d/\sa/\sb/\sc/\sd/\tag/\pe/\cv in {
  0/{1/2}/0/0/{1/2}/hot/cell/cell/hot/{(a) 同向配对}/{1}/{1/4},
  42/0/{1/2}/{1/2}/0/cell/hot/hot/cell/{(b) 反向配对}/{0}/{-1/4},
  84/{1/4}/{1/4}/{1/4}/{1/4}/hot/hot/hot/hot/{(c) 独立配对}/{1/2}/{0}}{
 \begin{scope}[xshift=\dx mm]
  \node[FigureInk] at (4.5mm,6.5mm) {$y=0$};
  \node[FigureInk] at (13.5mm,6.5mm) {$y=1$};
  \node[FigureInk,anchor=east] at (-1mm,0) {$x=0$};
  \node[FigureInk,anchor=east] at (-1mm,-6.5mm) {$x=1$};
  \node[\sa] at (4.5mm,0) {$\a$};
  \node[\sb] at (13.5mm,0) {$\b$};
  \node[\sc] at (4.5mm,-6.5mm) {$\c$};
  \node[\sd] at (13.5mm,-6.5mm) {$\d$};
  \node[marg] at (23mm,0) {$1/2$};
  \node[marg] at (23mm,-6.5mm) {$1/2$};
  \node[marg,minimum width=9mm] at (4.5mm,-14mm) {$1/2$};
  \node[marg,minimum width=9mm] at (13.5mm,-14mm) {$1/2$};
  \node[FigureMuted,font=\fontsize{7.5}{9}\selectfont,anchor=west] at (-6mm,-21.5mm) {$\mathbb P(X=Y)=\pe$，$\operatorname{Cov}=\cv$};
  \node[FigureInk] at (9mm,-28mm) {\tag};
 \end{scope}}
\end{tikzpicture}
